\chapter{结论与展望}\label{chap7_conclusion_outlook}

\section{主要研究结论}\label{sec_ch6_conclusion}

本文面向塔吊吊装作业中的近场安全防护需求，构建了涵盖硬件部署、多传感器标定、点云融合定位、动态目标检测、视觉语义识别及安全预警的完整研究框架。与既有研究中侧重单一图像检测或距离监测的思路不同，本文着重提升吊钩端近场空间的全方位感知能力，并将预警判断建立在三维空间几何关系与动态趋势分析之上。主要研究结论可概括为以下四个方面。

（1）完成了吊钩端异构感知系统的架构设计与多传感器联合标定，为后续算法研究提供了统一的观测空间。本文针对吊钩端空间有限、振动剧烈、遮挡复杂的工况特点，提出了双激光雷达与工业相机相结合的异构配置方案，通过侧向对称布置与俯仰角优化实现近场盲区互补与覆盖增强，并建立了跨模态数据关联的几何基础与多源数据时间一致性保障机制。

（2）构建了基于双雷达融合点云的吊钩端定位建图与动态目标检测方法，实现了平台自身位姿估计与近场动态环境的同步感知。研究表明，双雷达实时融合配合近场ROI裁剪能够在保证有效观测的同时显著降低后续处理计算负荷；LIO-SAM融合框架可在吊钩端摆动与回转工况下维持相对稳定的位姿估计与静态地图构建。动态检测与短时跟踪链路能够为风险研判提供动态目标的运动状态与趋势信息。

（3）构建了基于分割掩码与点云投影回查的吊装构件空间定位方法，打通了从二维语义到三维安全表达的关键环节。研究表明，将LiDAR点云逐点投影至图像平面并依据分割掩码回查目标点云集合，是一种兼顾几何可靠性与语义可解释性的融合路径，所提出的双基准安全包络建模策略为吊装场景中的安全边界计算提供了统一框架。

（4）在已有动态检测与轨迹预测基础上，形成了面向吊装作业的多层级预警与避障设计框架。现场实验验证了静态距离阈值判级与动态TTC协同升级机制的有效性，系统在相向接近工况下能够提供有效的预警提前量，且全链路时延满足实时性要求。研究表明，静态距离阈值卡控可作为风险筛查的基础层，动态点检测与轨迹预测为相向接近判断和TTC风险升级提供时间维度的依据，三级预警状态机与解耦的控制建议机制使系统研究范畴从感知可行性扩展到风险量化分级与动作建议承接的完整闭环。

总体而言，本文的实践表明，吊装安全预警系统的各个功能环节——从传感器标定、融合定位、动态检测到视觉识别与风险判级——并非彼此孤立的模块，而是围绕吊装构件这一核心对象串联而成的空间语义链路。这条链路的完整性决定了系统能否从基础感知延伸至预警输出，乃至最终的安全辅助控制。

\section{论文主要创新点}\label{sec_ch6_innovation}

本文工作的核心贡献在于，针对吊钩端这一高危区域，从传感器部署方式到信息融合路径都做了与既有研究不同的尝试，主要体现在以下两个方面。

（1）在硬件部署层面，本文将感知单元直接安装于吊钩下方而非传统的塔臂或塔身位置，并采用双LiDAR侧向对称布置来弥补单雷达的视野盲区。这一部署方式使传感器随吊物同步运动，从根本上缩短了传感距离。与既有研究侧重塔身或塔臂端宏观防碰撞不同，本文方案覆盖了从硬件部署、数据融合到定位跟踪的完整环节，将感知粒度推进至吊钩端近场空间，在感知架构上形成了明显差异。

（2）在预警策略层面，本文没有沿用现有的单一距离阈值或纯图像检测思路，而是让视觉和LiDAR各自发挥所长：视觉用YOLO26实例分割识别吊装构件，LiDAR提供高精度三维点云，通过点云投影与掩码回查将两者自然衔接，得到带有明确语义的构件三维位置和空间包络。基于这一融合结果，本文又区分了静态包络球和动态包络圆柱两种几何基准，前者用于评估与固定结构的净空，后者用于判断与移动目标的关系，并引入TTC和预测净空作为风险升级的依据。最终通过迟滞状态机输出稳定的预警等级，形成了一套从几何约束到趋势预判的分层决策机制。

\section{研究不足与未来展望}\label{sec_ch6_outlook}

本文在吊钩端感知系统的构建、融合定位与动态检测、构件识别与空间定位、预警与避障设计等方面均取得了阶段性成果，但仍有几个问题未能充分解决。

（1）多传感器标定的长期稳定性缺乏充分验证。本系统的外参标定依赖离线流程，但在吊钩端持续振动、温度变化和机械冲击作用下，已标定的内外参可能出现缓慢漂移，进而影响点云投影精度。目前尚未评估这种漂移在实际工况中的退化速率，也没有建立在线校正机制。后续需要针对不同振动烈度和持续时间进行长期观测，明确标定参数的漂移规律，并据此设计轻量化的在线自监测与补偿策略。

（2）LIO-SAM的定位与动态检测在极端运动工况下表现尚不充分。吊钩端的大幅摆动、急停急起以及点云稀疏区域的遮挡，都会干扰帧间配准和动态点判别的稳定性。当前的阈值和参数设定在面对这些工况时缺乏自适应能力，位姿连续性和动态召回率仍有下降风险。此外，当前的风险预测采用常速度模型外推吊装构件与动态目标的未来位置，在回转启停等钟摆式晃动工况下，该假设与实际运动模式存在偏差。尽管指数平滑已对速度振荡有所抑制且预测净空取保守最小值，但持续大幅摆动时外推精度的定量退化程度仍需进一步评估。后续工作应针对吊装作业特有的运动模式，重点解决参数的自适应调整问题，包括LIO-SAM的关键阈值、动态检测的判别门限以及聚类参数在不同场景下的敏感度分析，并探索可适应非线性运动的预测方法。

（3）多层级预警与相向运动风险升级已在真实数据上完成验证，详细实验见第6章。从结果来看，静态判级准确率92.48\%、动态判级准确率86.65\%表明判级逻辑基本成立，相向接近工况下0.70~s至1.60~s的提前量也说明TTC机制能够提供有效的时间裕度。但当前预警链路止步于控制接口原型，尚未接入真实塔吊的执行机构，距离完整的感知—判级—控制闭环还有差距。控制接口中的限速比例、保持条件、人工确认规则等参数也需要通过与现场实操数据的反复比对来标定。后续应优先完成控制接口与执行机构的联调验证，使预警输出真正能够驱动设备响应，而不仅停留在建议层面。

（4）实验验证主要集中于常规光照与气象条件，尚未系统评估视觉退化场景下的系统鲁棒性。施工现场存在夜间作业、强光直射、严重扬尘等典型工况，这些条件下视觉感知能力显著下降。在本系统的雷视融合架构中，吊装构件状态依赖于分割掩码约束下的点云投影回查，视觉退化将直接导致掩码缺失或边界不稳定，进而使跨模态投影关联中断，吊物的包络参数无法有效更新。此时系统仅能退守至LIO-SAM提供的吊钩端位姿作为预警基准，但吊物相对于吊钩的偏移与摆动状态难以全面感知。后续工作应针对上述退化场景开展专项实验，明确各模态失效的临界条件，并设计面向退化场景的降级感知策略与多模态互补机制，以维持系统在全天候工况下的基本预警能力。

% 参考文献请使用主文件中的 biblatex + \printbibliography 自动生成。
